\documentclass[11pt]{article}
\usepackage[a4paper,margin=1in]{geometry}
\usepackage{amsmath,amssymb,amsthm,microtype}
\usepackage[hidelinks]{hyperref}
\usepackage[nameinlink,noabbrev]{cleveref}
\usepackage[numbers]{natbib}
\newtheorem{theorem}{Theorem}[section]
\newtheorem{proposition}[theorem]{Proposition}
\newtheorem{corollary}[theorem]{Corollary}
\newtheorem{remark}[theorem]{Remark}
\newcommand{\Zorb}{\mathcal Z_{\rm orb}}
\newcommand{\Zsharp}{\mathcal Z^{\sharp}}
\newcommand{\Dtwo}{\mathbb D^2}
\newcommand{\Hm}{\mathcal H_m}
\newcommand{\Ps}{\Psi}
\hypersetup{pdftitle={The Two-Fugacity Divisor of the Weighted Henon Reflection Euler Family},pdfauthor={Liang Wang},pdfsubject={Weighted reflection Euler products and analytic hypersurface divisors},pdfkeywords={Henon map, reflection symmetry, dynamical zeta, Mobius inversion, analytic divisor, essential singularity}}
\title{The Two-Fugacity Divisor of the\linebreak Weighted H\'enon Reflection Euler Family}
\author{Liang Wang\\\small School of Artificial Intelligence and Automation\\\small Huazhong University of Science and Technology\\\small Wuhan 430074, P. R. China}
\date{August 16, 2026}

\begin{document}
\maketitle

\begin{abstract}
The orbit-resolved reflection Euler product of the full H\'enon horseshoe
depends on a positive defect weight $q$, but its previously known
all-channel continuation was restricted to $q=1$.  We lift each orbit
monomial by setting $w=qz$, so that
$z^nq^{S_n\chi}=z^{n-S_n\chi}w^{S_n\chi}$.  Primitive M\"obius subtraction
and Euler repetition then give the exact two-variable identity
\[
 \log\mathcal Z^{\sharp}(z,w)
 =\sum_{m\ge1}c_m\frac{2w^m}{1-z^{2m}-w^{2m}},\qquad
 c_m=\frac1m\prod_{\substack{p\mid m\\p\ \mathrm{odd}}}(1-p).
\]
Every $c_m$ is nonzero.  In the bidisk, the hypersurfaces
$\mathcal H_m=\{z^{2m}+w^{2m}=1\}$ are smooth and form a locally finite
effective divisor, while the channel sum converges normally on its
complement.  Restriction to $w=qz$, $q>0$, recovers the weighted orbit
product.  The $m$th channel has $2m$ explicitly located roots and a nonzero
closed principal coefficient at each root, so every one is exponentially
essential.  We do not infer a dense natural boundary, introduce a weighted
Lind zeta, or construct an operator or arithmetic trace.
\end{abstract}

\section{The weighted packet and the missing global parameter}

For the area-preserving map $H(x,y)=(6-x^2-y,x)$, the inherited horseshoe
coding identifies the odd reflection packets used below with
reversal-fixed binary words \citep{DevaneyNitecki1979,Arai2007}.  Let
$A_n$ be the primitive marked packet at odd period $n$ and define
\[
 \chi(s)=\mathbf 1\{s_{-1}=s_1\},\qquad
 S_n\chi(s)=\sum_{j=0}^{n-1}\chi(\sigma^js).
\]
The orbit-resolved product is
\begin{equation}\label{eq:original-product}
 \Zorb(z,q)=\prod_{\substack{n\ge1\\n\ \mathrm{odd}}}
 \prod_{\omega\in A_n}
 \left(1-z^nq^{S_n\chi(\omega)}\right)^{-1},\qquad q>0.
\end{equation}
It is a marked reflection-packet product, not the full Lind zeta of a flip
action.

The preceding transfer calculation gives, for $n=2h+1$,
\begin{align}
 F_n(q)&=2q(1+q^2)^h,\label{eq:all-polynomial}\\
 E_n(q)&=\sum_{k\mid n}\mu(k)F_{n/k}(q^k),\label{eq:primitive-polynomial}
\end{align}
where $F_n$ counts all reflection words and $E_n$ counts the primitive
marked packet with its individual energy weights.  If
$E(z,q)=\sum_{n\ \mathrm{odd}}E_n(q)z^n$, expansion of the Euler factors gives
\begin{equation}\label{eq:repetition}
 \log\Zorb(z,q)=\sum_{r\ge1}\frac1rE(z^r,q^r)
\end{equation}
in the disk of absolute convergence.  At $q=1$, grouping the M\"obius and
repetition indices exposes infinitely many rational channels.  Our goal is
to retain the full weight without postulating a new external zeta formula.

\section{A two-fugacity lift}

The inequality $0\le S_n\chi\le n$ permits a bivariate lift of every orbit
factor.  Set $w=qz$ and write
\begin{equation}\label{eq:monomial-lift}
 z^nq^{S_n\chi}=z^{n-S_n\chi}w^{S_n\chi}.
\end{equation}
We denote by $\Zsharp(z,w)$ the resulting formal Euler product.  Equation
\eqref{eq:all-polynomial} immediately becomes a rational all-word function.

\begin{proposition}[Lifted all-word function]\label{prop:all-word}
For sufficiently small $z,w$,
\begin{equation}\label{eq:Fsharp}
 F^{\sharp}(z,w)
 :=\sum_{h\ge0}2w(z^2+w^2)^h
 =\frac{2w}{1-z^2-w^2}.
\end{equation}
Primitive subtraction and repetition take the forms
\begin{align}
 E^{\sharp}(z,w)
 &=\sum_{\substack{k\ge1\\k\ \mathrm{odd}}}
   \mu(k)F^{\sharp}(z^k,w^k),\label{eq:Esharp}\\
 \log\Zsharp(z,w)
 &=\sum_{r\ge1}\frac1rE^{\sharp}(z^r,w^r).\label{eq:Zsharp-repeat}
\end{align}
On $w=qz$, these identities agree coefficientwise with
\eqref{eq:primitive-polynomial}--\eqref{eq:repetition}.
\end{proposition}

\begin{proof}
In a reflection word of length $2h+1$, the fixed boundary edge contributes
one unit of energy.  Each of the remaining $h$ paired edges either changes
the symbol, contributing $z^2$, or preserves it, contributing $w^2$; the
initial binary symbol contributes the factor $2w$.  This proves
\eqref{eq:Fsharp}.  A period-$n/k$ word repeated $k$ times dilates both
exponents in \eqref{eq:monomial-lift} by $k$, so ordinary odd M\"obius
inversion gives \eqref{eq:Esharp}.  Expanding $-\log(1-u)$ yields
\eqref{eq:Zsharp-repeat}.  Substitution $w=qz$ reverses
\eqref{eq:monomial-lift} and recovers the original formulas.
\end{proof}

\section{Exact channel regrouping}

For $m\ge1$, define
\begin{equation}\label{eq:Psi}
 \Ps_m(z,w)=\frac{2w^m}{1-z^{2m}-w^{2m}}.
\end{equation}

\begin{theorem}[Weighted scalar-channel identity]\label{thm:channels}
In the initial domain of absolute convergence,
\begin{equation}\label{eq:channels}
 \log\Zsharp(z,w)=\sum_{m\ge1}c_m\Ps_m(z,w),
\end{equation}
where
\begin{equation}\label{eq:cm}
 c_m=\frac1m\sum_{\substack{k\mid m\\k\ \mathrm{odd}}}k\mu(k)
 =\frac1m\prod_{\substack{p\mid m\\p\ \mathrm{odd}}}(1-p).
\end{equation}
For every $m$, $c_m\ne0$ and $|c_m|\le1$.  Consequently, for fixed $q>0$,
\begin{equation}\label{eq:qchannels}
 \log\Zorb(z,q)=\sum_{m\ge1}c_m
 \frac{2(qz)^m}{1-(1+q^{2m})z^{2m}}.
\end{equation}
\end{theorem}

\begin{proof}
Insert \eqref{eq:Esharp} into \eqref{eq:Zsharp-repeat}.  A pair consisting
of an odd M\"obius index $k$ and a repetition index $r$ contributes
\[
 \frac{\mu(k)}r\frac{2w^{kr}}
 {1-z^{2kr}-w^{2kr}}.
\]
Grouping by $m=kr$ gives the divisor sum in \eqref{eq:cm}.  That sum factors
over the odd radical of $m$, proving the product formula.  No factor $1-p$
vanishes.  Moreover
\[
 |c_m|\le \frac1m\prod_{\substack{p\mid m\\p\ \mathrm{odd}}}p\le1.
\]
Finally, $w=qz$ in \eqref{eq:channels} gives \eqref{eq:qchannels}.
\end{proof}

The coefficient $c_m$ is independent of $q$: all weight dependence has
moved into the geometry of the channel denominator.  This separation is the
main reason for using two fugacities rather than treating $q$ as a passive
real parameter.

\section{The bidisk hypersurface divisor}

Let $\Dtwo=\{(z,w)\in\mathbb C^2:|z|<1,|w|<1\}$ and define
\begin{equation}\label{eq:Hm}
 \Hm=\{(z,w)\in\Dtwo:z^{2m}+w^{2m}=1\},\qquad
 \mathcal H=\bigcup_{m\ge1}\Hm.
\end{equation}

\begin{theorem}[Smooth locally finite divisor]\label{thm:divisor}
Every $\Hm$ is a smooth reduced hypersurface.  The formal sum
$\sum_{m\ge1}\Hm$ is a locally finite effective analytic divisor in the
bidisk.  The series in \eqref{eq:channels} converges normally on compact
subsets of $\Dtwo\setminus\mathcal H$ and therefore defines a holomorphic
function $\mathcal L^{\sharp}$ there.  Its exponential is a nonvanishing
scalar continuation of the initial Euler germ.
\end{theorem}

\begin{proof}
The gradient of $z^{2m}+w^{2m}-1$ is
\[
 (2mz^{2m-1},2mw^{2m-1}).
\]
It cannot vanish on $\Hm$, since simultaneous vanishing would imply
$z=w=0$.  Thus $\Hm$ is smooth and reduced.

For a compact $K\subset\Dtwo$, choose $r,s<1$ such that $|z|\le r$ and
$|w|\le s$ on $K$.  Once $m$ is large enough,
$r^{2m}+s^{2m}<1$, so $K\cap\Hm$ is empty.  This proves local finiteness.

Now suppose $K\subset\Dtwo\setminus\mathcal H$.  The finitely many early
denominators in \eqref{eq:Psi} have a positive minimum modulus on $K$.
For all sufficiently large $m$,
\[
 |1-z^{2m}-w^{2m}|\ge1-r^{2m}-s^{2m}\ge\tfrac12.
\]
Together with $|c_m|\le1$, this gives
$|c_m\Ps_m(z,w)|\le4s^m$.  The Weierstrass test proves normal convergence.
Agreement with the initial logarithm follows from \cref{thm:channels}, and
exponentiation gives the asserted continuation.
\end{proof}

The theorem does not classify intersections $\Hm\cap\mathcal H_j$.  Such
joint complex collisions are unnecessary for the fixed-positive-weight
fibers considered next.

\section{Moving roots on positive-weight fibers}

Fix $q>0$ and restrict to $w=qz$.  The $m$th denominator in
\eqref{eq:qchannels} has the roots
\begin{equation}\label{eq:roots}
 \alpha_{m,\ell}(q)=\rho_m(q)e^{\pi i\ell/m},\qquad
 \rho_m(q)=(1+q^{2m})^{-1/(2m)},\qquad 0\le\ell<2m.
\end{equation}

\begin{proposition}[Strict radius separation]\label{prop:radii}
For every $q>0$, $\rho_m(q)$ is strictly increasing in $m$ and
\begin{equation}\label{eq:radius-limit}
 \lim_{m\to\infty}\rho_m(q)=\min(1,q^{-1}).
\end{equation}
Hence two different channels have no common root on a fixed positive-$q$
fiber.
\end{proposition}

\begin{proof}
It suffices to show that
$x\mapsto x^{-1}\log(1+q^{2x})$ is strictly decreasing.  Put
$y=q^{2x}$.  The numerator controlling its derivative is the negative of
\[
 (1+y)\log(1+y)-y\log y,
\]
which is positive for $y>0$.  This proves strictness.  The limit follows by
separating $q<1$, $q=1$, and $q>1$ in the defining formula.  Distinct radii
preclude a common complex root.
\end{proof}

\begin{theorem}[Exact fiberwise principal coefficient]\label{thm:principal}
Let $q>0$, $m\ge1$, and $0\le\ell<2m$.  Near
$\alpha=\alpha_{m,\ell}(q)$,
\begin{equation}\label{eq:principal}
 \log\Zorb(z,q)=
 \frac{c_m(-1)^\ell q^m}{m\sqrt{1+q^{2m}}}
 \frac1{1-z/\alpha}+G_{m,\ell,q}(z),
\end{equation}
where $G_{m,\ell,q}$ is holomorphic.  Thus the scalar continuation has an
exponential essential singularity at every root in \eqref{eq:roots}.
\end{theorem}

\begin{proof}
Set $v=1-z/\alpha$.  Since
$\alpha^{2m}=(1+q^{2m})^{-1}$ and
$\alpha^m=(-1)^\ell/\sqrt{1+q^{2m}}$,
\begin{align*}
 1-(1+q^{2m})z^{2m}&=2mv+O(v^2),\\
 2(qz)^m&=\frac{2(-1)^\ell q^m}{\sqrt{1+q^{2m}}}+O(v).
\end{align*}
By \cref{prop:radii}, no other channel denominator vanishes at $\alpha$.
Normal convergence from \cref{thm:divisor} makes their sum holomorphic in a
small neighborhood.  Multiplication by the nonzero $c_m$ gives
\eqref{eq:principal}; exponentiating a nonzero simple pole gives an
essential singularity.
\end{proof}

The radii and their limit locate the divisor but do not by themselves prove
that complex roots approach every point of the limiting circle.  We reserve
that angular accumulation and the resulting natural-boundary question for a
separate theorem.

\section{Executable certificate and claim boundary}

The primary certificate compares \eqref{eq:qchannels}, coefficient by
coefficient as exact polynomials in $q$, with the primitive/repetition law
through degree $48$.  An independent sparse-polynomial implementation
extends the comparison through degree $64$ without importing the primary
module.  The test suite checks degree $100$, reconstructs the first $128$
channel coefficients, and audits all complex roots through channel $24$ at
$q=1/2,1,2$.  Dependency hashes lock the proofs, certificates, and PDFs of
P69, P70, and P72.  These finite checks guard the implementation; the
all-$m$ statements follow from the proofs above.

The resulting object has a genuine two-variable analytic divisor and exact
fiberwise singular data.  We make no comparison with a Lind zeta away from
$q=1$, because no such source formula has been supplied.  We construct no
transfer, nuclear, trace-class, or self-adjoint operator, and the integer
channel label $m$ has no certified rational-prime or von-Mangoldt meaning.
No arithmetic advance or Route-B authorization follows.

\bibliographystyle{plainnat}
\bibliography{references}
\end{document}
